\documentclass[red]{beamer}
\usetheme{Warsaw}

% Specify character encoding
\usepackage[T1]{fontenc}
\usepackage[utf8]{inputenc}

% Add support for images
\usepackage{graphicx}

% Polish characters
\usepackage{fontspec}
\usepackage[polish]{babel}

\title{GHRS: Graph-based Hybrid Recommendation System with Application to Movie Recommendation}
\author[ K. Kostka, R Jarzyński, M. Jeżewski, M. Gagoś ]{
    Kamil Kostka \and
    Radosław Jarzyński \and
    Maciej Jeżewski \and
    Michał Gagoś
}
\date{Lublin, 7 listopada 2026}

\begin{document}

    % Display authors
    \begin{frame}
    \titlepage
    \end{frame}

    \section{Artykuł}
    \subsection{Dane artykułu}
    \begin{frame}{Artykuł bazowy}
        \begin{itemize}
            \item \textbf{Tytuł:} GHRS: Graph-based Hybrid Recommendation System with Application to Movie Recommendation
            \item \textbf{Autorzy:} Zahra Zamanzadeh Darban (Scopus h-index: 4), Mohammad Hadi Valipour (Scopus h-index: 4)
            \item \textbf{Czasopismo:} Expert Systems with Applications
            \item \textbf{Wydawca:} Elsevier
            \item \textbf{IF:} 9.4, \textbf{punkty MEiN:} 200
            \item \textbf{Źródło:} Preprint w repozytorium arXiv (\url{https://arxiv.org/abs/2111.11293})
        \end{itemize}
    \end{frame}

    \subsection{Cel}
    \begin{frame}{Cel badań i słowa kluczowe}
        \small
        \begin{itemize}
            \item \textbf{Główny problem:} „zimny start” (cold start) -- niezdolność klasycznych systemów do precyzyjnego polecania filmów nowym użytkownikom bez historii ocen
            \item \textbf{Cel artykułu:} hybrydowy system rekomendacyjny (GHRS) łączący grafowe miary podobieństwa użytkowników z ich profilem demograficznym, aby zniwelować skutki braku początkowych danych
            \item \textbf{Kluczowa teza:} przepuszczenie surowych danych (demografii i statystyk grafowych) przez autoenkoder wyodrębnia najważniejsze, ukryte cechy; system filtruje szum i dokładniej dzieli użytkowników na grupy o zbieżnym guście
        \end{itemize}
        \begin{block}{Słowa kluczowe}
        systemy rekomendacyjne, uczenie głębokie, modelowanie oparte na grafach, autoenkoder, problem zimnego startu
        \end{block}
    \end{frame}

    \section{Badania}
    \subsection{Materiał}
    \begin{frame}{Materiał badawczy}
        Dwa rzeczywiste zbiory z platformy MovieLens, wybrane ze względu na dane demograficzne (wiek, płeć, zawód, kod pocztowy).
        \medskip
        \begin{center}
            \begin{tabular}{lrrr}
                \textbf{Zbiór} & \textbf{Oceny (1--5)} & \textbf{Użytkownicy} & \textbf{Filmy}\\ \hline
                MovieLens 100K & 100\,000 & 943 & 1\,682\\
                MovieLens 1M & 1\,000\,209 & 6\,040 & ok.\ 3\,900
            \end{tabular}
        \end{center}
    \end{frame}

    \subsection{Metody}
    \begin{frame}{Metody badawcze i implementacja}
        \small
        \begin{enumerate}
            \item \textbf{Graf podobieństw:} węzły to użytkownicy, krawędzie łączą użytkowników o podobnym odsetku wspólnych ocen; próg $\alpha = 0{,}005$--$0{,}03$
            \item \textbf{Cechy grafu:} PageRank, Closeness Centrality, Betweenness
            \item \textbf{Fuzja cech:} konkatenacja z danymi demograficznymi (one-hot encoding)
            \item \textbf{Redukcja wymiarów:} 5-warstwowy autoenkoder, ReLU, optymalizator Adam
            \item \textbf{Klastrowanie:} K-Means, $K$ metodą łokcia ($K=8$ dla 100K, $K=9$ dla 1M)
            \item \textbf{Estymacja ocen:} średnia preferencji członków klastra
        \end{enumerate}
    \end{frame}

    \section{Wyniki i nowość}
    \subsection{Wyniki}
    \begin{frame}{Wyniki badań}
        \small
        \begin{itemize}
            \item \textbf{Skuteczność główna:} na zbiorze 1M GHRS osiągnął RMSE $=0{,}833$ (precyzja $0{,}792$, czułość $0{,}838$); na zbiorze 100K RMSE $=0{,}887$
            \item \textbf{Porównanie z konkurencją:} na 1M GHRS ma najniższy RMSE w zestawieniu, m.in. wobec SVD++ ($0{,}856$), Autoencoder COFILS ($0{,}838$) i RBM-CF ($0{,}872$); na 100K wynik równy AutoRec
            \item \textbf{Zimny start:} gdy ukryto oceny do 20\% użytkowników, błąd rósł w kontrolowany sposób (1M: z ok.\ $0{,}83$ do ok.\ $0{,}96$) -- model radzi sobie z rekomendacjami na podstawie samej demografii
        \end{itemize}
    \end{frame}

    \subsection{Nowość}
    \begin{frame}{Elementy nowości}
        \footnotesize
        \begin{itemize}
            \item \textbf{Fuzja relacji grafowych z demografią:} topologia grafu (jak użytkownicy łączą się przez podobne oceny) i profile demograficzne w jednej, wspólnej macierzy
            \item \textbf{Inna rola autoenkodera:} „kompresor” i filtr -- wyciąga z połączonej macierzy najważniejsze wzorce (cechy o niskiej korelacji), odrzucając szum
            \item \textbf{Klastrowanie na „czystych” danych:} K-Means działa na skompresowanych cechach z autoenkodera, nie na danych wejściowych; klastry są dokładniejsze, a generowanie rekomendacji szybsze
        \end{itemize}
        \begin{alertblock}{Ograniczenie modelu}
            Zależność od danych pobocznych (demografii) -- brak zastosowania na platformach zbierających tylko login (bez płci i wieku).
        \end{alertblock}
    \end{frame}

    \section{Wkład}
    \subsection{Odtwarzalność}
    \begin{frame}{Odtwarzalność badań i dane źródłowe}
        \begin{itemize}
            \item \textbf{Odtwarzalność:} badania są w pełni odtwarzalne, eksperyment zachowuje pełną transparentność metodologiczną
            \item \textbf{Dane:} zbiory MovieLens to standard w dziedzinie systemów rekomendacyjnych, darmowe i otwarte
            \item \textbf{Kod źródłowy:} \url{https://github.com/hadoov/GHRS}
        \end{itemize}
    \end{frame}

    \subsection{Wkład własny}
    \begin{frame}{Wkład własny}
        \footnotesize
        \textbf{Tytuł:} „Ewaluacja miękkich metod klastrowania (GMM, FCM) w grafowo-hybrydowych systemach rekomendacyjnych zmagających się z problemem zimnego startu”
        \begin{itemize}
            \item \textbf{Miękkie klastrowanie:} GMM i Fuzzy C-Means zamiast K-Means -- prawdopodobieństwa przynależności do wielu klastrów, gładsza estymacja dla użytkowników o zróżnicowanych gustach
            \item \textbf{Wzbogacenie metryki:} krawędzie grafu uzupełnione o kosinusowe podobieństwo gatunków (TF-IDF z metadanych filmów) obok progu $\alpha$
            \item \textbf{Nowa architektura:} VAE zamiast autoenkodera (ciągła przestrzeń ukryta ułatwia zimny start) oraz optymalizator Nadam
            \item \textbf{Wdrożenie:} Python (PyTorch, scikit-learn)
            \item \textbf{Efekty:} mniejsza wariancja rekomendacji, lepsza odporność na wartości odstające
        \end{itemize}
    \end{frame}

\end{document}
